\SourceLeaf{8}
\pdfbookmark[0]{Untersuchung über die elliptischen Elemente der Pallas}{palladis}
\begin{center}
{\large UNTERSUCHUNG}\par\medskip
{\Large ÜBER DIE ELLIPTISCHEN ELEMENTE DER PALLAS}\par\medskip
AUS DEN OPPOSITIONEN DER JAHRE 1803, 1804, 1805, 1807, 1808, 1809.
\end{center}
\Article{1}
Der im März 1802 entdeckte Hauptplanet Pallas ist bisher siebenmal in periodischer Folge erschienen. Da der Planet aber am Tag seiner Entdeckung die Opposition zur Sonne schon überschritten hatte, zählt man bis jetzt nur sechs Oppositionen; die meisten davon wurden an mehreren Sternwarten recht gut beobachtet und konnten daher mit ausreichender Genauigkeit bestimmt werden. Ausnehmen muss man allein die Opposition von 1808. In diesem Jahr befand sich der Planet im Aphel seiner Bahn und bereitete den Beobachtern wegen seines schwachen Lichts erhebliche Schwierigkeiten. Um die Wahrheit zu sagen, haben ihn die Astronomen ein wenig vernachlässigt. Der angesehene \textsc{von Lindenau} lieferte zwar eine recht große Zahl von Rektaszensionen, die er um die Zeit der Opposition mit gewohnter Sorgfalt am ausgezeichneten Meridianfernrohr der Seeberger Sternwarte beobachtet hatte. Von anderer Seite konnte ich jedoch nur wenige und weit weniger sichere Deklinationen erhalten. Deshalb ließ sich diese Opposition nicht mit derselben Genauigkeit wie die übrigen bestimmen, und besonders die Breite blieb zu unsicher.
\Article{2}
Sobald die Beobachtungen einer Erscheinung abgeschlossen waren, pflegte ich sie jedes Jahr vollständig zu sammeln, sorgfältig zu prüfen und mit den Beobachtungen der Vorjahre zu verbinden, um daraus die elliptischen Elemente umso \SourceLeaf{9}genauer bestimmen zu können. Aus der letzten Rechnung dieser Art, die ich gegen Ende des Jahres 1807 ausführte, gingen die Elemente hervor, die im Januarheft 1808 des XVII. Bandes des Commercium literarium astronomicum des angesehenen \textsc{von Zach} veröffentlicht wurden. Sie stimmten mit allen vier bis dahin beobachteten Oppositionen sehr gut überein. Im folgenden Jahr hielt mich nach der fünften Opposition der Mangel an zuverlässigen Beobachtungen von neuen Rechnungen ab. Ich nahm sie erst wieder auf, als mir die 1809 um die sechste Opposition angestellten Beobachtungen erlaubten, die Rechnung für beide Oppositionen gemeinsam in Angriff zu nehmen. Ich habe sie in unseren Gelehrten Anzeigen vom 24. Februar 1810 angekündigt.

\Needspace{11\baselineskip}
Hier folgt eine Übersicht über alle sechs bisher beobachteten Oppositionen:
\begin{center}\small\setlength{\tabcolsep}{3pt}
\begin{tabular}{l r r|r|r|r}
\multicolumn{3}{c|}{\shortstack{Oppositionszeit am\\Göttinger Meridian}}&\shortstack{Tage seit Beginn\\des Jahres 1803}&\shortstack{Heliozentrische\\Länge}&\shortstack{Geozentrische\\Breite}\\\hline
1803 Juni&30.&$0^{\mathrm h}\;27^{\mathrm m}\;32^{\mathrm s}$&181.019120&$277^0\;39'\;24''0$&$+46^0\;26'\;36''0$\\
1804 Aug.&30.&$4\;58\;27$&608.207257&$337\;0\;36.1$&$+15\;1\;49.8$\\
1805 Nov.&29.&$11\;15\;4$&1064.468796&$67\;20\;42.9$&$-54\;30\;54.9$\\
1807 Mai&4.&$14\;37\;41$&1585.609502&$223\;37\;27.7$&$+42\;11\;25.6$\\
1808 Juli&26.&$21\;17\;32$&2034.887176&$304\;2\;59.7$&$+37\;43\;53.7$\\
1809 Sept.&22.&$16\;10\;20$&2457.673843&$359\;40\;4.4$&$-7\;22\;10.1$
\end{tabular}\end{center}
\Article{3}
Wenn Pallas sich genau auf einer Ellipse und nach den Gesetzen \textsc{Kepler}s bewegte, könnten die neuesten Elemente, die mit den vier Oppositionen der Jahre 1803, 1804, 1805, 1807 sehr gut übereinstimmen, nur noch sehr kleine Fehler enthalten. Sie müssten gewiss auch mit den nächsten Oppositionen bis auf eine oder zwei Bogenminuten übereinstimmen. Tatsächlich war die Übereinstimmung so weit davon entfernt, dass diese Elemente schon bei der fünften Opposition eine Abweichung von bis zu vier Bogenminuten zeigten und bei der sechsten um zwölf Bogenminuten abwichen. Unser Planet erfährt nämlich durch die übrigen Planeten, besonders durch Jupiter, so große Störungen, dass sich zwischen den Erscheinungen und einer rein elliptischen Bewegung keine genaue und dauerhafte Übereinstimmung erzielen lässt. Zugleich ist daraus ersichtlich, dass sich je nach den vier zugrunde gelegten Oppositionen jeweils andere Elemente ergeben. Das hat die Berechnung zweier weiterer Elementensysteme bestätigt, von denen ich das eine kürzlich aus den Oppositionen \SourceLeaf{10}der Jahre 1804, 1805, 1807, 1808 und das andere aus denen der Jahre 1805, 1807, 1808, 1809 abgeleitet habe. Damit sich diese verschiedenen Elementensysteme besser vergleichen lassen, gebe ich sie hier einzeln an.

\Needspace{14\baselineskip}
\begin{center}I.\quad\textit{Elliptische Elemente der Pallas aus den Oppositionen der Jahre}\par
1803, 1804, 1805, 1807.\end{center}
\begingroup\small\setlength{\tabcolsep}{2pt}
\noindent\begin{tabular}{@{}p{.51\linewidth}lr@{}}
Epoche der mittleren Länge für 1803 am Göttinger Meridian&. . . . . . . . . . . . . .&$221^0\;39'\;30''4$\\
Mittlere tägliche tropische Bewegung&. . . . . . . . . . . .&$770''2143$\\
Länge des Perihels 1803&. . . . . . . . . . . . .&$121^0\;3'\;11''4$\\
Länge des aufsteigenden Knotens 1803&. . . . . . . . . . .&$172\;28\;56.9$\\
Bahnneigung&. . . . . . . . . . . . . . . .&$34\;37\;41.0$\\
Exzentrizität $(=\sin.\,14^0\;10'\;58''81)$&. . . . . . . . . .&$0.2450198$\\
Logarithmus der großen Halbachse&. . . . . . . . . . . . .&$0.4423149$
\end{tabular}\par\endgroup

\Needspace{13\baselineskip}
\begin{center}II.\quad\textit{Elliptische Elemente der Pallas aus den Oppositionen der Jahre}\par
1804, 1805, 1807, 1808.\end{center}
\begingroup\small\setlength{\tabcolsep}{2pt}
\noindent\begin{tabular}{@{}p{.51\linewidth}lr@{}}
Epoche der mittleren Länge 1803&. . . . . . . . . . .&$221^0\;34'\;56''7$\\
Mittlere tropische tägliche Bewegung&. . . . . . . . . . . .&$770''4467$\\
Länge des Perihels 1803&. . . . . . . . . . . . .&$121^0\;5'\;22''1$\\
Länge des aufsteigenden Knotens 1803&. . . . . . . . . . .&$172\;28\;46.8$\\
Bahnneigung&. . . . . . . . . . . . . . . .&$34\;37\;31.5$\\
Exzentrizität $(=\sin.\,14^0\;10'\;4''08)$&. . . . . . . . . . .&$0.2447624$\\
Logarithmus der großen Halbachse&. . . . . . . . . . . . .&$0.4422276$
\end{tabular}\par\endgroup

\Needspace{13\baselineskip}
\begin{center}III.\quad\textit{Elliptische Elemente der Pallas aus den Oppositionen der Jahre}\par
1805, 1807, 1808, 1809.\end{center}
\begingroup\small\setlength{\tabcolsep}{2pt}
\noindent\begin{tabular}{@{}p{.51\linewidth}lr@{}}
Epoche der mittleren Länge 1803&. . . . . . . . . . .&$221^0\;23'\;24''6$\\
Mittlere tropische tägliche Bewegung&. . . . . . . . . . . .&$770''9265$\\
Länge des Perihels 1803&. . . . . . . . . . . . .&$120^0\;58'\;4''8$\\
Länge des aufsteigenden Knotens 1803&. . . . . . . . . . .&$172\;27\;52.4$\\
Bahnneigung&. . . . . . . . . . . . . . . .&$34\;36\;49.4$\\
Exzentrizität $(=\sin.\,14^0\;9'\;36''63)$&. . . . . . . . . . .&$0.2446335$\\
Logarithmus der großen Halbachse&. . . . . . . . . . . . .&$0.4420473$
\end{tabular}\par\endgroup
\SourceLeaf{11}
\Article{4}
Wie sich die elliptische Bahn eines Planeten aus vier Beobachtungen bestimmen lässt, von denen nur zwei vollständig sind, habe ich in Abschnitt II von Buch II der Theoria motus corporum coelestium gezeigt. Die dort angegebene Lösung behandelt die Aufgabe in größter Allgemeinheit. Im Sonderfall, dass die vier Beobachtungen Oppositionen sind, empfiehlt sich jedoch ein anderes Verfahren. Ich hoffe, dass seine Darstellung den Astronomen willkommen sein wird. Dieses Verfahren setzt zwar voraus, dass man jedes der Elemente näherungsweise kennt, oder wenigstens die vier Elemente Bahnneigung, Länge des Knotens, Länge des Perihels und Exzentrizität. Es spricht aber nichts dagegen, diese Kenntnis aus einer früheren Rechnung als bereits vorhanden vorauszusetzen. Das Wesentliche hängt von der Lösung der folgenden Aufgabe ab.
\begin{quote}\itshape
„Zu vier gegebenen Längen des Planeten in seiner Bahn, die gegebenen Zeitpunkten entsprechen, die Länge des Perihels, die Exzentrizität, die mittlere tägliche Bewegung und die Epoche der mittleren Länge finden“
\end{quote}
Die Lösung gebe ich im folgenden Artikel an.
\Article{5}
Da die Länge des Perihels als bereits näherungsweise bekannt vorausgesetzt wird, erhält man durch Abziehen dieser Länge von den Längen in der Bahn die näherungsweisen wahren Anomalien $v$, $v'$, $v''$, $v'''$. Dabei kann man jene Länge entweder als konstant ansehen oder, wenn ihre jährliche Änderung bekannt ist, auf die einzelnen Zeitpunkte entsprechend übertragen. Sei $e=\sin\varphi$ die näherungsweise Exzentrizität. Mit ihr berechne man nach den bekannten Verfahren die mittleren Anomalien $M$, $M'$, $M''$, $M'''$, die den wahren Anomalien $v$, $v'$, $v''$, $v'''$ entsprechen. Bezeichnen wir nun die Zeiten mit $t$, $t'$, $t''$, $t'''$, dann müssten ihre Differenzen $t'-t$, $t''-t'$, $t'''-t''$ den Differenzen $M'-M$, $M''-M'$, $M'''-M''$ proportional sein, wenn wir von den wahren Werten der Länge des Perihels und des Winkels $\varphi$ ausgegangen wären. Andernfalls erhält man die an diesen Größen anzubringenden Korrekturen auf folgende Weise. Setzt man
\[
\begin{aligned}
m&=-\frac{\cos\varphi^3}{(1+e\cos v)^2},&n&=\frac{m(2+e\cos v)\sin v}{\cos\varphi'}\\[6pt]
m'&=-\frac{\cos\varphi^3}{(1+e\cos v')^2},&n'&=\frac{m'(2+e\cos v')\sin v'}{\cos\varphi}\\[6pt]
m''&=-\frac{\cos\varphi^3}{(1+e\cos v'')^2},&n''&=\frac{m''(2+e\cos v'')\sin v''}{\cos\varphi}\\[6pt]
m'''&=-\frac{\cos\varphi^3}{(1+e\cos v''')^2},&n'''&=\frac{m'''(2+e\cos v''')\sin v'''}{\cos\varphi}
\end{aligned}
\]
\noindent\SourceLeaf{12}
so folgt aus Artikel 15 der Theoria Motus Corporum Coelestium, dass sich bei einer Korrektur des Perihels um den kleinen Zuwachs $\mathrm d\Pi$ und des Winkels $\varphi$ um den Zuwachs $\mathrm d\varphi$ für die mittleren Anomalien die korrigierten Werte ergeben
\[
\begin{gathered}
M+m\mathrm d\Pi+n\mathrm d\varphi\\
M'+m'\mathrm d\Pi+n'\mathrm d\varphi\\
M''+m''\mathrm d\Pi+n''\mathrm d\varphi\\
M'''+m'''\mathrm d\Pi+n'''\mathrm d\varphi
\end{gathered}
\]
sofern man Potenzen und Produkte der Korrekturen $\mathrm d\Pi$, $\mathrm d\varphi$ vernachlässigen darf. Man erhält also drei Ausdrücke für die mittlere tägliche Bewegung in Bezug auf das Perihel, nämlich
\[
\begin{gathered}
\frac{M'-M}{t'-t}+\frac{m'-m}{t'-t}\mathrm d\Pi+\frac{n'-n}{t'-t}\mathrm d\varphi\\[5pt]
\frac{M''-M'}{t''-t'}+\frac{m''-m'}{t''-t'}\mathrm d\Pi+\frac{n''-n'}{t''-t'}\mathrm d\varphi\\[5pt]
\frac{M'''-M''}{t'''-t''}+\frac{m'''-m''}{t'''-t''}\mathrm d\Pi+\frac{n'''-n''}{t'''-t''}\mathrm d\varphi
\end{gathered}
\]
Aus ihrer Gleichheit ergeben sich die Unbekannten $\mathrm d\Pi$, $\mathrm d\varphi$. Setzt man diese in einen der Ausdrücke ein, so ergibt sich die korrigierte mittlere Bewegung, aus der nach dem bekannten Verfahren auch die große Halbachse bestimmt wird. Setzt man schließlich die Werte der Korrekturen $\mathrm d\Pi$, $\mathrm d\varphi$ in den Ausdruck für eine beliebige mittlere Anomalie ein, so ergibt sich daraus ihr korrigierter Wert und damit die auf einen beliebigen Zeitpunkt zu übertragende Epoche der mittleren Länge.

Übrigens ist leicht einzusehen, wie sich bei Bedarf sowohl die periodischen als auch die säkularen Störungen berücksichtigen lassen. Man muss die gegebenen Längen in der Bahn lediglich von den periodischen Störungen befreien, bevor man sie in die Rechnung einführt. Mithilfe der säkularen Störungen sind dagegen die für eine beliebige Epoche angenommenen Werte für die Länge des Perihels und die Exzentrizität auf die einzelnen Zeitpunkte $t$, $t'$, $t''$, $t'''$ zu übertragen.
\Article{6}
Sehen wir nun, wie sich dies zur Bestimmung der Bahn aus vier beobachteten Oppositionen verwenden lässt. Bei der ersten Opposition sei die Länge des Planeten in der Ekliptik $\alpha$, die geozentrische Breite $\beta$, der Abstand der Erde von der Sonne $R$, die näherungsweise Länge des aufsteigenden Knotens $\AscendingNode$ und die näherungsweise Neigung der Bahnebene gegen die Ekliptik $i$. Die zu den übrigen Oppositionen gehörenden Größen bezeichne ich mit entsprechenden Zeichen, die sich durch einen, zwei oder drei Indizes unterscheiden. Aus $\alpha$, $\AscendingNode$, $i$ ergibt sich die Länge in der Bahn nach der bekannten
\SourceLeaf{13}Formel. Den drei übrigen Längen in der Ekliptik legt man eine entsprechende Rechnung zugrunde. Aus diesen vier Längen in der Bahn leitet man nach den Vorschriften des vorhergehenden Artikels die elliptischen Elemente ab und daraus nach den bekannten Formeln die vier Radiusvektoren $r$, $r'$, $r''$, $r'''$. Setzt man nun
\[
\frac{R\sin\beta}{r}=\sin(\beta-\gamma),\qquad \frac{R'\sin\beta'}{r'}=\sin(\beta'-\gamma')\quad\text{etc.}
\]
so sind $\gamma$, $\gamma'$, $\gamma''$, $\gamma'''$ offenbar die vier heliozentrischen Längen, die aus den geozentrischen Breiten abgeleitet wurden. Leiten wir dieselben Größen aus den Längen ab, indem wir setzen
\[
\operatorname{tang}i\sin(\alpha-\AscendingNode)=\operatorname{tang}\delta,\qquad \operatorname{tang}i\sin(\alpha'-\AscendingNode)=\operatorname{tang}\delta'\ \text{etc.}
\]
Offenbar müsste dann $\gamma=\delta$, $\gamma'=\delta'$, $\gamma''=\delta''$, $\gamma'''=\delta'''$ gelten, wenn für $\AscendingNode$ und $i$ die wahren Werte eingesetzt worden wären und sich alle vier Oppositionen, durch die die Aufgabe überbestimmt ist, mit derselben Ellipse darstellen ließen. Ergibt sich etwas anderes, so erhält man die Korrekturen $\mathrm d\AscendingNode$, $\mathrm di$ für die Größen $\AscendingNode$, $i$ auf folgende Weise. Soweit diese Korrekturen als kleine Größen erster Ordnung gelten, deren Potenzen und Produkte man vernachlässigen darf, ergeben sich bei Verwendung der Werte $\AscendingNode+\mathrm d\AscendingNode$, $i+\mathrm di$ anstelle der Breiten $\gamma$, $\delta$, $\gamma'$, $\delta'$ usw. offenbar andere Breiten von folgender Form
\[
\begin{array}{ll}
\gamma+a\mathrm d\AscendingNode+b\mathrm di,&\delta+c\mathrm d\AscendingNode+f\mathrm di\\
\gamma'+a'\mathrm d\AscendingNode+b'\mathrm di,&\delta'+c'\mathrm d\AscendingNode+f'\mathrm di\quad\text{etc.}
\end{array}
\]
wobei man die Koeffizienten $c$, $f$, $c'$, $f'$ usw. nach den folgenden Formeln findet
\[
\begin{array}{ll}
c=-\frac12\sin2\delta\operatorname{cotang}(\alpha-\AscendingNode),&f=\dfrac{\sin2\delta}{\sin2i}\\[5pt]
c'=-\frac12\sin2\delta'\operatorname{cotang}(\alpha'-\AscendingNode),&f'=\dfrac{\sin2\delta'}{\sin2i}\quad\text{etc.}
\end{array}
\]
Die Koeffizienten $a$, $b$ usw. könnten wir ebenfalls auf analytischem Weg bestimmen, doch ziehe ich das folgende Verfahren vor. Ich wiederhole die eben beschriebene Rechnung und setze anstelle der Knotenlänge $\AscendingNode$ einen anderen Wert ein, der sich von ihr um einen kleinen, beliebig gewählten Betrag unterscheidet. Daraus lassen sich die Zahlenwerte der Koeffizienten $a$, $a'$ usw. leicht gewinnen. Wiederholt man dieselbe Rechnung mit einer geringfügig geänderten Bahnneigung, erhält man entsprechend die Werte der Koeffizienten $b$, $b'$ usw. Zugleich zeigt sich so, wie groß die Korrekturen der elliptischen Elemente ausfallen, wenn die Größen $\AscendingNode$, $i$ kleine Korrekturen erfahren.

\SourceLeaf{14}
Damit liegen zur Bestimmung der Korrekturen $\mathrm d\AscendingNode$, $\mathrm di$ \textit{vier} Gleichungen vor
\[
\begin{gathered}
\gamma-\delta+(a-c)\mathrm d\AscendingNode+(b-f)\mathrm di=0\\
\gamma'-\delta'+(a'-c')\mathrm d\AscendingNode+(b'-f')\mathrm di=0\\
\gamma''-\delta''+(a''-c'')\mathrm d\AscendingNode+(b''-f'')\mathrm di=0\\
\gamma'''-\delta'''+(a'''-c''')\mathrm d\AscendingNode+(b'''-f''')\mathrm di=0
\end{gathered}
\]
Wegen der unvermeidlichen Beobachtungsfehler lassen sich diese Gleichungen kaum genau erfüllen. Die am besten geeigneten Werte der Unbekannten $\mathrm d\AscendingNode$, $\mathrm di$ werden daher nach den in Abschnitt III von Buch II der Theoria Motus Corporum Coelestium erläuterten Grundsätzen bestimmt. Daraus lassen sich nach dem eben Gesagten zugleich die Korrekturen der übrigen Elemente leicht ableiten,
\Article{7}
Diesen Vorschriften möchte ich noch einige Bemerkungen hinzufügen.

I. Es bedarf kaum des Hinweises, dass unsere Formeln auch dann anwendbar sind, wenn die Knoten nicht als ruhend betrachtet werden. Man braucht den angenommenen Wert der Knotenlänge nur von seiner Epoche auf die vier Zeitpunkte umzurechnen und diese vier verschiedenen Werte in den Formeln für die jeweiligen vier Örter zu verwenden. Dasselbe gilt für die Bahnneigung, sofern ihre säkulare Änderung bekannt geworden ist und ihre Berücksichtigung lohnend erscheint. Will man außerdem die Störungen der heliozentrischen Breite in die Rechnung aufnehmen, so muss man sie in den vorhergehenden Gleichungen offenbar entweder von $\gamma$, $\gamma'$, $\gamma''$, $\gamma'''$ abziehen oder zu $\delta$, $\delta'$, $\delta''$, $\delta'''$ addieren.

III. Wenn die Bahnneigung und die Exzentrizität mäßig sind, werden die Koeffizienten $a$, $b$, $a'$, $b'$ usw. so klein, dass man sie vernachlässigen und damit auf die oben vorgeschriebene zweite und dritte Rechnung verzichten kann. Hat man dann die Korrekturen $\mathrm d\AscendingNode$, $\mathrm di$ gefunden, so muss man, sofern sie nicht zufällig sehr klein ausfallen, die Längen in der Bahn aus den korrigierten Werten der Knotenlänge und der Bahnneigung, also $\AscendingNode+\mathrm d\AscendingNode$, $i+\mathrm di$, neu berechnen. Weichen die korrigierten Längen in der Bahn um mehrere Bogensekunden von den früheren unkorrigierten ab, so muss man die elliptischen Elemente nach Artikel 5 erneut berechnen. Übrigens wird man die Koeffizienten $m$, $n$, $m'$ usw. kaum je neu berechnen müssen, da sich ihre Werte schon aus der ersten Annahme hinreichend genau ergeben.
\SourceLeaf{15}
\Article{8}
Zur besseren Erläuterung der vorhergehenden Vorschriften gebe ich hier die vollständige Rechnung an, mit der das dritte Elementensystem aus den Oppositionen der Jahre 1805, 1807, 1808, 1809 bestimmt wurde. Ich setzte die Länge des aufsteigenden Knotens für den Beginn des Jahres 1803 zu $172^0\,28'\,46''8=\AscendingNode$ und die Bahnneigung zu $34^0\,37'\,31''5$ an. Jene Länge rechnete ich auf die Zeitpunkte der einzelnen Oppositionen um, indem ich die Präzession $2'\,26''01$, $3'\,37''50$, $4'\,39''12$, $5'\,37''11$ addierte. Anschließend leitete ich aus den Längen die heliozentrischen Breiten ab
\[
\begin{array}{rrr}
\delta=-33^0&40'&50''63\\
\delta'=+28&14&51.24\\
\delta''=+27&20&55.86\\
\delta'''=-4&52&28.44
\end{array}
\]
\[
\begin{array}{ll}
c=+0.1252,&f=-0.9870\\
c'=-0.3366.&f'=+0.8917\\
c''=+0.3609.&f''=+0.8727\\
c'''=+0.6803,&f'''=-0{,}1811
\end{array}
\]
Ferner ergeben sich die Breitenargumente
\[
\begin{array}{rrr}
257^0&25'&6''73\\
56&24&5.96\\
126&2&55.07\\
188&36&2.69
\end{array}
\]
aus denen die elliptischen Elemente zu ermitteln sind. Dazu setze ich den Winkel $\varphi=14^0\,10'\,4''08$, wie er im zweiten Elementensystem gefunden wurde, und die Länge des Perihels für den Beginn des Jahres 1803 $=121^0\,5'\,22''1$. Das Perihel wird ebenso wie der Knoten gegenüber den Fixsternen als unbeweglich vorausgesetzt. Sein Abstand vom aufsteigenden Knoten ist daher konstant $=308^0\,36'\,35''3$. Nach den bekannten Formeln erhalten wir damit die folgenden Werte.
\[
\begin{array}{rrr@{\quad}rrr}
v=308^0&48'&31''43,&M=328^0&15'&45''08\\
v'=107&47&30.66,&M'=79&46&27.05\\
v''=177&26&19.77,&M''=175&54&28.87\\
v'''=239&59&27.39,&M'''=266&29&57.59
\end{array}
\]
\SourceLeaf{16}
\[
\begin{array}{ll}
m=-0.68517,&n=+1.18569\\
m'=-1.06482,&n=-2.01318\\
m''=-1.59701,&n=-0.12921\\
m'''=-1.18352,&n=+1.93947
\end{array}
\]
Bezeichnen wir also die Korrektur der Länge des Perihels mit $\mathrm d\Pi$ und die Korrektur des Winkels $\varphi$ mit $\mathrm d\varphi$, so ergibt sich die mittlere siderische Bewegung

von der ersten bis zur zweiten Opposition innerhalb von $521.140706$ Tagen
\[
111^0\,30'\,41''97-0.37965\mathrm d\Pi-3.19887\mathrm d\varphi
\]
von der zweiten bis zur dritten Opposition innerhalb von $449.277674$ Tagen
\[
96^0\,8'\,1''82-0.53219\mathrm d\Pi+1.88397\mathrm d\varphi
\]
von der dritten bis zur vierten Opposition innerhalb von $422.786667$ Tagen
\[
90^0\,35'\,28''72+0.41349\mathrm d\Pi+2.06868\mathrm d\varphi
\]
Daraus ergeben sich drei Ausdrücke für die mittlere tägliche siderische Bewegung
\[
\begin{gathered}
770''31398-0.0007285\mathrm d\Pi-0.0061382\mathrm d\varphi\\
770.30718-0.0011845\mathrm d\Pi+0.0041933\mathrm d\varphi\\
771.37892+0.0009780\mathrm d\Pi+0.0048930\mathrm d\varphi
\end{gathered}
\]
aus deren Vergleich sich die folgenden Gleichungen ergeben
\[
\begin{gathered}
0=0''00680+0.0004560\mathrm d\Pi-0.0103315\mathrm d\varphi\\
0=-1.07174-0.0021625\mathrm d\Pi-0.0006997\mathrm d\varphi
\end{gathered}
\]
daraus $\mathrm d\Pi=-488''82$, $\mathrm d\varphi=-20''92$, und folglich
\begin{center}\fontsize{9}{11}\selectfont\setlength{\tabcolsep}{3pt}\begin{tabular}{lr}
korrigierte Länge des Perihels zu Beginn des Jahres 1803,&$120^0\,57'\,13''28$,\\
und korrigierter Wert des Winkels $\ .\ .\ .\ .\ .\ .$&$\varphi=14^0\,9'\,43''16$\\
ferner ergibt sich die mittlere tägliche siderische Bewegung $\ .\ .\ .\ .\ .$&$770''7985$\\
korrigierte mittlere Anomalie für die erste Opposition $\ .$&$328^0\,20'\,55''20$\\
mittlere Länge für dieselbe Epoche $\ .\ .\ .\ .\ .$&$89^0\,20'\,34''49$
\end{tabular}\end{center}
Schließlich ergibt sich aus der mittleren siderischen Bewegung der Logarithmus der großen Halbachse $0.4420439$ und damit der Logarithmus des Halbparameters $0.4152361$. Da wir nun die korrigierten Werte der wahren Anomalien haben
\SourceLeaf{17}
\[
\begin{array}{rrr}
v=308^0&56'&40''25\\
v'=107&55&39.48\\
v''=177&34&28.59\\
v'''=240&7&36.21
\end{array}
\]
berechnet man die Logarithmen der Radiusvektoren
\[
\begin{aligned}
\log r&=0.3531088\\
\log r'&=0.4492406\\
\log r''&=0.5369700\\
\log r'''&=0.4716739
\end{aligned}
\]
Aus den Sonnentafeln haben wir ferner
\[
\begin{aligned}
\log R&=9.9937332\\
\log R'&=0.0039862\\
\log R''&=0.0065917\\
\log R'''&=0.0011160
\end{aligned}
\]
woraus wir schließlich ableiten
\[
\begin{array}{rrr}
\gamma=-33^0&39'&48''15\\
\gamma'=+28&15&0.73\\
\gamma''=+27&20&9.07\\
\gamma'''=-4&52&53.99
\end{array}
\]

Um die Koeffizienten $a$, $b$ usw. zu erhalten, bilde ich eine \textit{zweite Annahme}. Ich behalte die Bahnneigung bei, vergrößere aber die Knotenlänge um eine Bogenminute, sodass sie für den Beginn des Jahres 1803 $172^0\,29'\,46''8$ beträgt. Setzt man dann wie in der ersten Annahme die Länge des Perihels für dieselbe Epoche $=121^0\,5'\,22''1$, so findet man die wahren Anomalien
\[
\begin{array}{rrr}
v=308^0&48'&40''93\\
v'=107&47&34.06\\
v''=177&26&22.25\\
v'''=239&59&15.00
\end{array}
\]
und daraus, wenn man wieder $\varphi=14^0\,10'\,4''08$ setzt, die mittleren Anomalien
\SourceLeaf{18}
\[
\begin{array}{rrr}
M=328^0&15'&51''59\\
M'=79&46&30.67\\
M''=175&54&32.83\\
M'''=266&29&42.93
\end{array}
\]
(Für diese Rechnung ist das übliche Verfahren übrigens entbehrlich. Es genügt, zu den in der ersten Annahme gefundenen Werten der mittleren Anomalien die Produkte aus den positiv genommenen Koeffizienten $m$, $m'$, $m''$, $m'''$ und den Differenzen zwischen den entsprechenden Werten der wahren Anomalien in beiden Annahmen zu addieren, nämlich
\[
\begin{aligned}
+9''50\times0.68517&=+6''51\\
+3.40\times1.06482&=+3.62\\
+2.48\times1.59701&=+3.96\\
-12.39\times1.18352&=-14.66)
\end{aligned}
\]
Behält man die Werte der Koeffizienten $m$, $n$ usw. bei, so erhält man daraus
\begin{center}\fontsize{9}{11}\selectfont\setlength{\tabcolsep}{3pt}\begin{tabular}{p{.23\linewidth}p{.70\linewidth}}
$\mathrm d\Pi=-468''21$,&Länge des Perihels zu Beginn des Jahres 1803, $120^0\,57'\,33''89$\\
$\mathrm d\varphi=-20''62$,&$\varphi=14^0\,9'\,43''46$
\end{tabular}\end{center}
\begin{center}\fontsize{9}{11}\selectfont\setlength{\tabcolsep}{3pt}\begin{tabular}{ll}
mittlere tägliche siderische Bewegung&$=770''7761$\\
Logarithmus der großen Halbachse&$=0.4420523$\\
mittlere Länge bei der ersten Opposition&$=89^0\,20'\,47''85$
\end{tabular}\end{center}
\[
\begin{array}{rrr}
\gamma=-33^0&39'&51''10\\
\gamma'=+28&15&1.27\\
\gamma''=+27&20&10.97\\
\gamma'''=-4&52&54.66
\end{array}
\]

Schließlich bilde ich eine \textit{dritte Annahme}, indem ich $\AscendingNode=172^0\,28'\,46''6$, $i=34^0\,38'\,31''5$ setze. Daraus ergibt sich wie zuvor
\begin{center}\fontsize{9}{11}\selectfont\setlength{\tabcolsep}{3pt}\begin{tabular}{lr}
Länge des Perihels für den Beginn des Jahres 1803,&$120^0\,55'\,34''46$\\
&$\varphi=14^0\,9'\,52''63$\\
mittlere tägliche siderische Bewegung&$=770''8398$\\
Logarithmus der großen Halbachse&$=0.4420283$\\
mittlere Länge bei der ersten Opposition&$=89^0\,20'\,20''65$
\end{tabular}\end{center}
\SourceLeaf{19}
\[
\begin{array}{rrr}
\gamma=-33^0&39'&35''63\\
\gamma'=+28&15&5.20\\
\gamma''=+27&20&9.32\\
\gamma'''=-4&52&52.65
\end{array}
\]

Der Vergleich der Werte der heliozentrischen Breiten $\gamma$, $\gamma'$, $\gamma''$, $\gamma'''$, die in den drei Annahmen gefunden wurden, liefert
\[
\begin{array}{ll}
a=-0.0492,&b=+0.2087\\
a'=+0.0212,&b'=+0.0745\\
a''=+0.0317,&b''=+0.0042\\
a'''=-0.0112,&b'''=+0.0223
\end{array}
\]
Man erhält somit vier Gleichungen
\[
\begin{gathered}
+62''48-0.1744\mathrm d\AscendingNode+1.1957\mathrm di=0\\
+9.49+0.3578\mathrm d\AscendingNode-0.8172\mathrm di=0\\
-46.79-0.3292\mathrm d\AscendingNode-0.8685\mathrm di=0\\
-25.55-0.6915\mathrm d\AscendingNode+0.2034\mathrm di=0
\end{gathered}
\]
aus denen sich nach dem Prinzip der kleinsten Quadrate ergibt
\[
\begin{aligned}
\mathrm d\AscendingNode&=-54''41\\
\mathrm di&=-42.06
\end{aligned}
\]
sodass man erhält
\begin{center}\fontsize{9}{11}\selectfont\setlength{\tabcolsep}{3pt}\begin{tabular}{lr}
Länge des aufsteigenden Knotens für den Beginn des Jahres 1803,&$172^0\,27'\,52''39$\\
Bahnneigung $\ .\ .\ .\ .\ .\ .\ .\ .\ .\ .$&$34\quad36\quad49.44$
\end{tabular}\end{center}

Die übrigen Elemente lassen sich entweder durch Vergleich der Ergebnisse der drei Annahmen gewinnen oder, genauer, durch eine neue Berechnung der Längen in der Bahn und durch Wiederholung der in Artikel 5 erläuterten Rechenschritte. Das erste Verfahren liefert
\begin{center}\fontsize{9}{11}\selectfont\setlength{\tabcolsep}{3pt}\begin{tabular}{lr}
Länge des Perihels 1803 $\ .\ .\ .\ .\ .\ .\ .$&$120^0\,58'\,3''86$\\
Winkel $\varphi\ .\ .\ .\ .\ .\ .\ .\ .\ .\ .\ .\ .\ .$&$14\quad9\quad36.25$\\
mittlere Länge bei der ersten Opposition $\ .\ .$&$89\quad20\quad32.08$\\
mittlere tägliche siderische Bewegung $\ .\ .\ .\ .\ .\ .\ .$&$770''7899$\\
Logarithmus der großen Halbachse. $\ .\ .\ .\ .\ ,\ .\ .$&$0.4420471$
\end{tabular}\end{center}
\SourceLeaf{20}
Mit dem anderen Verfahren erhält man
\begin{center}\fontsize{9}{11}\selectfont\setlength{\tabcolsep}{3pt}\begin{tabular}{lr}
Länge des Perihels 1803 $\ .\ .\ .\ .\ .\ .\ .\ .$&$120^0\,58'\,4''81$\\
Winkel $\varphi\ .\ .\ .\ .\ .\ .\ .\ .\ .\ .\ .\ .\ .$&$14\quad9\quad36.63$\\
mittlere Länge bei der ersten Opposition $\ .\ .\ .\ .$&$89\quad20\quad31.81$\\
mittlere tägliche siderische Bewegung $\ .\ .\ .\ .\ .\ .\ .$&$770''7893$\\
Logarithmus der großen Halbachse $\ .\ .\ .\ .\ .\ .\ .\ .$&$0.4420473$
\end{tabular}\end{center}
Mit diesen Elementen stimmen die oben in Artikel 3 angegebenen überein.
\Article{9}
So groß die Störungen auch sein mögen, die Pallas durch die übrigen Planeten erfährt, so zeigt doch die Erfahrung, dass an vier Oppositionen angepasste elliptische Elemente der Bewegung des Planeten während dieses ganzen Zeitraums recht gut genügen. Sie weichen sogar von seiner vorhergehenden und nachfolgenden Bewegung nur wenig ab, sofern man den Zeitraum nicht zu groß wählt. So wich etwa das zweite der in Artikel 3 angegebenen Elementensysteme bei der Opposition von 1803 um drei und bei der Opposition von 1809 um fünf Bogenminuten von der beobachteten heliozentrischen Länge ab. Unter diesen Umständen scheint es mir zur Aufstellung einer Ephemeride für die künftige Bewegung des Planeten stets am zweckmäßigsten, rein elliptische Elemente zu verwenden, die aus den vier unmittelbar vorhergehenden Oppositionen abgeleitet wurden. Die Zahl der aus den Störungen entstehenden Gleichungen wird gewiss so groß werden, dass selbst die Berechnung eines einzigen heliozentrischen Orts des Planeten beinahe so aufwendig sein muss wie die Berechnung der elliptischen Elemente nach den oben angegebenen Vorschriften. So lässt sich der geozentrische Ort des Planeten stets mit einem Fehler von wenigen Bogenminuten sicher vorhersagen. Diese Genauigkeit genügt, um den Planeten aufzufinden.
\Article{10}
Doch die Würde der Wissenschaft verlangt, dass man auf eine dauerhaftere Übereinstimmung hinwirkt, die sich offensichtlich erst erzielen lässt, wenn die Störungen in die Rechnung aufgenommen werden. Nach meinem Urteil wäre eine so langwierige und mühsame Rechnung verfrüht gewesen, solange die verfügbaren Beobachtungen einen zu kurzen Zeitraum umfassten und sich die Störungen durch die übrigen Planeten kaum bemerkbar machten. Jetzt aber, da die elliptische Bewegung nicht mehr ausreicht, um alle beobachteten Örter miteinander in Einklang zu bringen, scheint die Zeit gekommen, in der man an eine genauere Theorie denken kann. Wie ich die Berechnung der Störungen, die Pallas vor allem durch Jupiter erfährt,
\SourceLeaf{21}
am zweckmäßigsten und genauesten ausführen zu können glaube, werde ich bald an anderer Stelle ausführlicher erläutern, da sich die bei anderen Planeten verwendeten Methoden wegen der übergroßen Exzentrizität und Bahnneigung kaum oder gar nicht anwenden lassen. Im Folgenden gehe ich noch auf die elliptischen Elemente ein, die sich als Grundlage der Störungsrechnung am besten zu eignen scheinen. Ich nehme mir vor, elliptische Elemente zu gewinnen, die nicht dieser oder jener Opposition genau genügen, sondern allen bisher beobachteten so gut wie möglich. Ein Verfahren zur Ausführung dieser Aufgabe habe ich bereits in Artikel 187 der Theoria Motus Corporum Coelestium kurz beschrieben. Was ich dort allgemein behandelt habe, lässt einige Abkürzungen zu, wenn die beobachteten Örter Oppositionen sind. Außerdem fehlen in jenem Werk einige praktische Kunstgriffe, mit denen ich seit Langem die Anwendung der Methode der kleinsten Quadrate zu erleichtern pflege. Ich hoffe daher, dass die Astronomen eine etwas ausführlichere Darstellung dieser Rechnungen hier begrüßen werden. Die ganze Aufgabe besteht darin, die \textit{Korrekturen} zu bestimmen, die zu den Näherungswerten der Elemente hinzuzufügen sind, von denen vorausgesetzt wird, dass sie von allen beobachteten Örtern nur wenig abweichen. Die ganze Arbeit hängt somit von zwei Schritten ab. Zunächst sind die linearen Gleichungen aufzustellen, die die einzelnen beobachteten Örter liefern. Dann sind aus diesen Gleichungen die am besten geeigneten Werte der Unbekannten zu gewinnen.
\Article{11}
Für die näherungsweisen Elemente bezeichnen die Symbole die folgenden Größen.
\begin{description}\setlength{\itemsep}{0pt}\setlength{\parsep}{0pt}
\item[$L$] die mittlere Länge des Planeten für eine beliebige Epoche
\item[$t$] die Zahl der Tage, die seit der Epoche bis zum Zeitpunkt der Beobachtung vergangen sind
\item[$\MeanMotionGlyph$] die mittlere tägliche siderische Bewegung in Bogensekunden
\item[$\Pi$] die Länge des Perihels
\item[$e=\sin\varphi$] die Exzentrizität
\item[$a$] die große Halbachse
\item[$r$] der Radiusvektor
\item[$v$] die wahre Anomalie
\item[$E$] die exzentrische Anomalie
\item[$\AscendingNode$] die Länge des aufsteigenden Knotens
\item[$i$] die Bahnneigung
\item[$u$] das Breitenargument
\item[$\lambda$] die heliozentrische Länge
\SourceLeaf{22}
\item[$\gamma$] die heliozentrische Breite
\item[$\ApproxBeta$] die geozentrische Breite
\item[$R$] der Abstand der Erde von der Sonne.
\end{description}
Aufgrund der Beobachtung setze ich dagegen die folgenden Werte an.
\begin{description}\setlength{\itemsep}{0pt}\setlength{\parsep}{0pt}
\item[$\alpha$] die heliozentrische Länge
\item[$\ObservedBeta$] die geozentrische Breite.
\end{description}
Schließlich bezeichne ich mit $\mathrm dL$, $\mathrm d\MeanMotionGlyph$, $\mathrm d\Pi$ usw. die Korrekturen der Größen $L$, $\MeanMotionGlyph$, $\Pi$ usw. Damit ist $\mathrm dL+t\mathrm d\MeanMotionGlyph$ die Korrektur der mittleren Länge,
$\mathrm dL+t\mathrm d\MeanMotionGlyph-\mathrm d\Pi$ ist die Korrektur der mittleren Anomalie,
und somit gilt nach Artikel 15 und 16 der Theoria Motus Corp. Coel.
\[
\begin{aligned}
\mathrm dv&=\frac{aa\cos\varphi}{rr}(\mathrm dL+t\mathrm d\MeanMotionGlyph-\mathrm d\Pi)+\frac{aa}{rr}(2-e\cos E-ee)\sin E\mathrm d\varphi\\
\mathrm dr&=\frac ra\mathrm da+a\operatorname{tang}\varphi\sin v(\mathrm dL+t\mathrm d\MeanMotionGlyph-\mathrm d\Pi)-a\cos\varphi\cos v\mathrm d\varphi
\end{aligned}
\]
Ferner ergibt sich als Korrektur des Breitenarguments $\mathrm du=\mathrm dv+\mathrm d\Pi-\mathrm d\AscendingNode$,
und nach Artikel 52 der Theoria Motus C. C. ergibt sich die Korrektur der heliozentrischen Länge aus
\[
\mathrm d\lambda=\mathrm d\AscendingNode-\operatorname{tang}\gamma\cos(\lambda-\AscendingNode)\mathrm di+\frac{\cos i}{\cos\gamma^2}\mathrm du
\]
Daraus folgt
\[
\begin{aligned}
\mathrm d\lambda={}&\frac{aa\cos\varphi\cos i}{rr\cos\gamma^2}\mathrm dL\\
&+\frac{taa\cos\varphi\cos i}{rr\cos\gamma^2}\mathrm d\MeanMotionGlyph\\
&+\left(\frac{\cos i}{\cos\gamma^2}-\frac{aa\cos\varphi\cos i}{rr\cos\gamma^2}\right)\mathrm d\Pi\\
&+\frac{aa\cos i}{rr\cos\gamma^2}(2-e\cos E-ee)\sin E\mathrm d\varphi\\
&+\left(1-\frac{\cos i}{\cos\gamma^2}\right)\mathrm d\AscendingNode\\
&-\operatorname{tang}\gamma\cos(\lambda-\AscendingNode)\mathrm di
\end{aligned}
\]
Da ferner gilt
\[
\begin{aligned}
a^{\frac32}\MeanMotionGlyph&=\mathrm{Const.}\\
r\sin(\ApproxBeta-\gamma)&=R\sin\ApproxBeta\\
\operatorname{tang}\gamma&=\operatorname{tang}i\sin(\alpha-\AscendingNode)
\end{aligned}
\]
ergibt sich durch Differentiation
\SourceLeaf{23}
\[
\frac{\mathrm da}{a}=-\frac{2\mathrm d\MeanMotionGlyph}{3\MeanMotionGlyph}
\]
\[
\frac{\mathrm dr}{r}+\operatorname{cotang}(\ObservedBeta-\gamma).(\mathrm d\ApproxBeta-\mathrm d\gamma)=\operatorname{cotg}\ApproxBeta\mathrm d\ApproxBeta,\quad\text{oder}
\]
\[
\mathrm d\ApproxBeta=\frac{\sin\ApproxBeta\cos(\ApproxBeta-\gamma)}{\sin\gamma}\mathrm d\gamma-\frac{\sin\ApproxBeta\sin(\ApproxBeta-\gamma)}{r\sin\gamma}\mathrm dr
\]
\[
\mathrm d\gamma=\frac{\sin2\gamma}{\sin2i}\mathrm di-\tfrac12\sin2\gamma\operatorname{cotang}(\alpha-\AscendingNode)\mathrm d\AscendingNode
\]
woraus sich mit dem oben entwickelten Ausdruck für $\mathrm dr$ ergibt
\[
\begin{aligned}
\mathrm d\ApproxBeta={}&-\frac{a\sin\ApproxBeta\sin(\ApproxBeta-\gamma)\operatorname{tang}\varphi\sin v}{r\sin\gamma}\mathrm dL\\
&+\left\{\frac{2\sin\ApproxBeta\sin(\ApproxBeta-\gamma)}{3\MeanMotionGlyph\sin\gamma}-\frac{at\sin\ApproxBeta\sin(\ApproxBeta-\gamma)\operatorname{tang}\varphi\sin v}{r\sin\gamma}\right\}\mathrm d\MeanMotionGlyph\\
&+\frac{a\sin\ApproxBeta\sin(\ApproxBeta-\gamma)\operatorname{tang}\varphi\sin v}{r\sin\gamma}\mathrm d\Pi\\
&+\frac{a\sin\ApproxBeta\sin(\ApproxBeta-\gamma)\cos\varphi\cos v}{r\sin\gamma}\mathrm d\varphi\\
&+\frac{2\sin\ApproxBeta\cos(\ApproxBeta-\gamma)\cos\gamma}{\sin2i}\mathrm di\\
&-\sin\ApproxBeta\cos(\ApproxBeta-\gamma)\cos\gamma\operatorname{cotang}(\alpha-\AscendingNode)\mathrm d\AscendingNode
\end{aligned}
\]
Als Werte der heliozentrischen Länge und der geozentrischen Breite ergeben sich aus den korrigierten Werten der Elemente also $\lambda+\mathrm d\lambda$, $\ApproxBeta+\mathrm d\ApproxBeta$. Damit liefert jede Opposition zwei Gleichungen
\[
\begin{aligned}
\alpha&=\lambda+\mathrm d\lambda\\
\ObservedBeta&=\ApproxBeta+\mathrm d\ApproxBeta
\end{aligned}
\]
\Article{12}
Wenden wir diese Vorschriften auf die sechs in Artikel 2 angegebenen Oppositionen der Pallas an
und legen wir der Rechnung das in Artikel 3 dargestellte zweite Elementensystem zugrunde, so erhalten wir die folgenden zwölf Gleichungen.

\noindent Aus der \textit{ersten} Opposition, mit der durch Rechnung gefundenen Länge $=277^0\,36'\,20''07$,
\[
\text{geozentrische Breite}=+46^0\,26'\,29''19:
\]
\[
\begin{aligned}
0={}&-183''93+0.79363\mathrm dL+143.66\mathrm d\MeanMotionGlyph+0.39493\mathrm d\Pi+0.95920\mathrm d\varphi\\
&-0.18856\mathrm d\AscendingNode+0.17387\mathrm di\\[4pt]
0={}&-6''81-0.02658\mathrm dL+46.71\mathrm d\MeanMotionGlyph+0.02658\mathrm d\Pi-0.20858\mathrm d\varphi\\
&+0.15946\mathrm d\AscendingNode+1.25782\mathrm di
\end{aligned}
\]
\SourceLeaf{24}
\Needspace{11\baselineskip}
Aus der \textit{zweiten} Opposition, mit der berechneten Länge $=337^0\,0'\,36''04$,
\[
\text{geozentrische Breite}=+15^0\,1'\,46''71:
\]
\[
\begin{aligned}
0={}&-0''06+0.58880\mathrm dL+358.12\mathrm d\MeanMotionGlyph+0.26208\mathrm d\Pi-0.85234\mathrm d\varphi\\
&+0.14912\mathrm d\AscendingNode+0.17775\mathrm di\\[4pt]
0={}&-3''09+0.01318\mathrm dL+28.39\mathrm d\MeanMotionGlyph-0.01318\mathrm d\Pi-0.07861\mathrm d\varphi\\
&+0.91704\mathrm d\AscendingNode+0.54365\mathrm di
\end{aligned}
\]
Aus der \textit{dritten} Opposition, mit der berechneten Länge $=67^0\,20'\,42''88$,
\[
\text{geozentrische Breite}=-54^0\,31'\,3''88:
\]
\[
\begin{aligned}
0={}&-0''02+1.73436\mathrm dL+1846.17\mathrm d\MeanMotionGlyph-0.54603\mathrm d\Pi-2.05662\mathrm d\varphi\\
&-0.18833\mathrm d\AscendingNode-0.17445\mathrm di\\[4pt]
0={}&-8''98-0.12606\mathrm dL-227.42\mathrm d\MeanMotionGlyph+0.12606\mathrm d\Pi-0.38939\mathrm d\varphi\\
&+0.17176\mathrm d\AscendingNode-1.35441\mathrm di
\end{aligned}
\]
Aus der \textit{vierten} Opposition, mit der berechneten Länge $=223^0\,37'\,25''39$,
\[
\text{geozentrische Breite}=+42^0\,11'\,28''07:
\]
\[
\begin{aligned}
0={}&-2''31+0.99584\mathrm dL+1579.03\mathrm d\MeanMotionGlyph+0.06456\mathrm d\Pi+1.99545\mathrm d\varphi\\
&-0.06040\mathrm d\AscendingNode-0.33750\mathrm di\\[4pt]
0={}&+2''47-0.08089\mathrm dL-67.22\mathrm d\MeanMotionGlyph+0.08089\mathrm d\Pi-0.09970\mathrm d\varphi\\
&-0.46359\mathrm d\AscendingNode+1.22803\mathrm di
\end{aligned}
\]
Aus der \textit{fünften} Opposition, mit der berechneten Länge $=304^0\,2'\,59''71$,
\[
\text{geozentrische Breite}=+37^0\,44'\,31''82:
\]
\[
\begin{aligned}
0={}&+0''01+0.65311\mathrm dL+1329.09\mathrm d\MeanMotionGlyph+0.38994\mathrm d\Pi-0.08439\mathrm d\varphi\\
&-0.04305\mathrm d\AscendingNode+0.34268\mathrm di\\[4pt]
0={}&+38''12-0.00218\mathrm dL+38.47\mathrm d\MeanMotionGlyph+0.00218\mathrm d\Pi-0.18710\mathrm d\varphi\\
&+0.47301\mathrm d\AscendingNode-1.14371\mathrm di
\end{aligned}
\]
Aus der \textit{sechsten} Opposition, mit der berechneten Länge $=359^0\,34'\,46''67$,
\[
\text{geozentrische Breite}=-7^0\,20'\,12''13:
\]
\[
\begin{aligned}
0={}&-317''73+0.69957\mathrm dL+1719.32\mathrm d\MeanMotionGlyph+0.12913\mathrm d\Pi-1.38787\mathrm d\varphi\\
&+0.17130\mathrm d\AscendingNode-0.08360\mathrm di\\[4pt]
0={}&+117''97-0.01315\mathrm dL-43.84\mathrm d\MeanMotionGlyph+0.01315\mathrm d\Pi+0.02929\mathrm d\varphi\\
&+1.02138\mathrm d\AscendingNode-0.27187\mathrm di
\end{aligned}
\]
Von diesen zwölf Gleichungen werden wir jedoch die zehnte ganz verwerfen, da die beobachtete geozentrische Breite zu unsicher ist.
\SourceLeaf{25}
\Article{13}
Da sich die sechs Unbekannten $\mathrm dL$, $\mathrm d\MeanMotionGlyph$ usw. nicht so bestimmen lassen, dass alle elf Gleichungen genau erfüllt sind, das heißt, dass die einzelnen Funktionen der Unbekannten auf der rechten Seite gleichzeitig $=0$ werden, werden wir diejenigen Werte ermitteln, für die die Summe der Quadrate dieser Funktionen möglichst klein ist. Sind allgemein die folgenden linearen Funktionen der Unbekannten $p$, $q$, $r$, $s$ usw. gegeben
\[
\begin{aligned}
n+ap+bq+cr+ds+\text{etc.}&=w\\
n'+a'p+b'q+c'r+d's+\text{etc.}&=w'\\
n''+a''p+b''q+c''r+d''s+\text{etc.}&=w''\\
n'''+a'''p+b'''q+c'''r+d'''s+\text{etc.}&=w'''
\end{aligned}
\]
usw., so sind offenbar die Bedingungsgleichungen dafür, dass $ww+w'w'+w''w''+w'''w'''+\text{etc.}=\Omega$ ein Minimum wird, die folgenden
\[
\begin{aligned}
aw+a'w'+a''w''+a'''w'''+\text{etc.}&=0\\
bw+b'w'+b''w''+b'''w'''+\text{etc.}&=0\\
cw+c'w'+c''w''+c'''w'''+\text{etc.}&=0\\
dw+d'w'+d''w''+d'''w'''+\text{etc.}&=0\\
\text{etc.\quad oder}&
\end{aligned}
\]
wenn wir der Kürze halber bezeichnen
\[
\begin{array}{ll}
an+a'n'+a''n''+a'''n'''+\text{etc.}&\text{mit }[an]\\
aa+a'a'+a''a''+a'''a'''+\text{etc.}&\text{mit }[aa]\\
ab+a'b'+a''b''+a'''b'''+\text{etc.}&\text{mit }[ab]\\
\text{etc.}&\\
bb+b'b'+b''b''+b'''b'''+\text{etc.}&\text{mit }[bb]\\
bc+b'c'+b''c''+b'''c'''+\text{etc.}&\text{mit }[bc]\\
\text{etc. etc.}&
\end{array}
\]
sind $p$, $q$, $r$, $s$ usw. durch Elimination aus den folgenden Gleichungen zu bestimmen
\[
\begin{aligned}{}
[an]+[aa]p+[ab]q+[ac]r+[ad]s+\text{etc.}&=0\\
[bn]+[ab]p+[bb]q+[bc]r+[bd]s+\text{etc.}&=0\\
[cn]+[ac]p+[bc]q+[cc]r+[cd]s+\text{etc.}&=0\\
[dn]+[ad]p+[bd]q+[cd]r+[dd]s+\text{etc.}&=0\\
\text{etc.}&
\end{aligned}
\]
\SourceLeaf{26}
Sobald die Zahl der Unbekannten $p$, $q$, $r$, $s$ usw. etwas größer ist, erfordert die Elimination jedoch eine sehr langwierige und mühsame Rechnung, die sich auf folgende Weise erheblich verkürzen lässt. Außer den Koeffizienten $[an]$, $[aa]$, $[ab]$ usw., deren Zahl $=\tfrac12(ii+3i)$ ist, wenn die Zahl der Unbekannten $=i$ ist, setze ich auch voraus, dass $nn+n'n'+n''n''+n'''n'''+\text{etc.}=[nn]$ berechnet worden ist. Dann ist leicht einzusehen, dass gilt
\[
\begin{aligned}
\Omega={}&[nn]+2[an]p+2[bn]q+2[cn]r+2[dn]s+\text{etc.}\\
&+[aa]pp+2[ab]pq+2[ac]pr+2[ad]ps+\text{etc.}\\
&\qquad+[bb]qq+2[bc]qr+2[bd]qs+\text{etc.}\\
&\qquad\qquad+[cc]rr+2[cd]rs+\text{etc.}\\
&\qquad\qquad\qquad+[dd]ss+\text{etc.}\\
&\qquad\qquad\qquad\qquad\text{etc.}
\end{aligned}
\]
Bezeichnen wir also
\[
[an]+[aa]p+[ab]q+[ac]r+[ad]s+\text{etc.}\quad\text{mit}\quad A
\]
so ist klar, dass alle Teile von $\dfrac{A^2}{[aa]}$, die den Faktor $p$ enthalten, in $\Omega$ vorkommen und somit $\Omega-\dfrac{A^2}{[aa]}$ eine von $p$ unabhängige Funktion ist. Setzen wir daher
\[
\begin{aligned}{}
[nn]-\frac{[an]^2}{[aa]}&=[nn,1]\\
[bn]-\frac{[an].[ab]}{[aa]}&=[bn,1]\\
[cn]-\frac{[an].[ac]}{[aa]}&=[cn,1]\\
[dn]-\frac{[an].[ad]}{[aa]}&=[dn,1]\quad\text{etc.}\\
[bb]-\frac{[ab]^2}{[aa]}&=[bb,1]\\
[bc]-\frac{[ab].[ac]}{[aa]}&=[bc,1]\\
[bd]-\frac{[ab].[ad]}{[aa]}&=[bd,1]\quad\text{etc. etc.: ergibt sich}
\end{aligned}
\]
\[
\begin{aligned}
\Omega-\frac{A^2}{[aa]}={}&[nn,1]+2[bn,1]q+2[cn,1]r+2[dn,1]s+\text{etc.}\\
&+[bb,1]qq+2[bc,1]qr+2[bd,1]qs+\text{etc.}\\
&\qquad+[cc,1]rr+2[cd,1]rs+\text{etc.}\\
&\qquad\qquad+[dd,1]ss+\text{etc.}\\
&\qquad\qquad\qquad\text{etc.}
\end{aligned}
\]
Diese Funktion bezeichnen wir mit $\Omega'$.

\SourceLeaf{27}
Setzen wir ebenso
\[
[bn,1]+[bb,1]q+[bc,1]r+[bd,1]s+\text{etc.}=B
\]
so ist $\Omega'-\dfrac{B^2}{[bb,1]}$ eine von $q$ unabhängige Funktion, die wir $=\Omega''$ setzen. Auf dieselbe Weise setzen wir
\[
\begin{aligned}{}
[nn,1]-\frac{[bn,1]^2}{[bb,1]}&=[nn,2]\\
[cn,1]-\frac{[bn,1].[bc,1]}{[bb,1]}&=[cn,2]\\
[cc,1]-\frac{[bc,1]^2}{[bb,1]}&=[cc,2]\\
\text{etc. etc.\quad und}&
\end{aligned}
\]
\[
[cn,2]+[cc,2]r+[cd,2]s+\text{etc.}=C
\]
sodass $\Omega''-\dfrac{C^2}{[cc,2]}$ auch von $r$ unabhängig ist. Wir fahren ebenso fort, bis wir in der Folge $\Omega$, $\Omega'$, $\Omega''$ usw. zu einem von allen Unbekannten unabhängigen Glied gelangen. Es lautet $[nn,\mu]$, wenn die Zahl der Unbekannten $p$, $q$, $r$, $s$ usw. mit $\mu$ bezeichnet wird. Wir haben also
\[
\Omega=\frac{A^2}{[aa]}+\frac{B^2}{[bb,1]}+\frac{C^2}{[cc,2]}+\frac{D^2}{[dd,3]}+\text{etc.}+[nn,\mu]
\]
Da nun $\Omega=ww+w'w'+w''w''\text{ etc.}$ seiner Natur nach keinen negativen Wert annehmen kann, lässt sich leicht zeigen, dass die Divisoren $[aa]$, $[bb,1]$, $[cc,2]$, $[dd,3]$ usw. notwendig \textit{positiv} sein müssen (der Kürze halber führe ich diesen Beweis hier jedoch nicht näher aus). Daraus folgt unmittelbar, dass sich der kleinste Wert von $\Omega$ ergibt, wenn man $A=0$, $B=0$, $C=0$, $D=0$ usw. setzt. Aus diesen $(\mu)$ Gleichungen sind also die Unbekannten $p$, $q$, $r$, $s$ usw. zu bestimmen. Das gelingt sehr leicht in umgekehrter Reihenfolge, da die letzte Gleichung offenbar nur eine Unbekannte enthält, die vorletzte zwei und so weiter. Zugleich hat dieses Verfahren den Vorzug, dass sich daraus der kleinste Wert der Summe $\Omega$ unmittelbar ergibt, da er offenbar $=[nn,\mu]$ ist.
\Article{14}
Wenden wir diese Vorschriften nun auf unser Beispiel an. Hier sind $p$, $q$, $r$, $s$ usw. die Größen $\mathrm dL$, $\mathrm d\MeanMotionGlyph$, $\mathrm d\Pi$, $\mathrm d\varphi$, $\mathrm d\AscendingNode$, $\mathrm di$. Nach sorgfältiger Durchführung der Rechnung fand ich die folgenden Zahlenwerte.
\SourceLeaf{28}
{\fontsize{9}{11}\selectfont
\[
\begin{array}{lll}
\begin{gathered}{}
[nn]=148848\\
[an]=-371.09\\
[bn]=-580104\\
[cn]=-113.45\\
[dn]=+268.53\\
[en]=+94.26\\
[fn]=-31.81\\
[aa]=+5.91569\\
[ab]=+7203.91\\
[ac]=-0.09344\\
[ad]=-2.28516\\
[ae]=-0.34664\\
[af]=-0.18194\\
[bb]=+10834225\\
[bc]=-49.06\\
[bd]=-3229.77\\
[be]=-198.64\\
[bf]=-143.05\\
[cc]=+0.71917\\
[cd]=+1.13382\\
[ce]=+0.06400\\
[cf]=+0.26341\\
[dd]=+12.00340\\
[de]=-0.37137\\
[df]=-0{,}11762\\
[ee]=+2.28215\\
[ef]=-0.36136\\
[ff]=+5.62456\\[3pt]
\text{Daraus leitete ich ferner ab}\\[3pt]
[nn,1]=+125569\\
[bn,1]=-138534
\end{gathered}&
\begin{gathered}{}
[cn,1]=-119.31\\
[dn,1]=-125.18\\
[en,1]=+72.52\\
[fn,1]=-43.22\\
[bb,1]=+2458225\\
[bc,1]=+62.13\\
[bd,1]=-510.58\\
[be,1]=+213.84\\
[bf,1]=+73.45\\
[cc,1]=+0.71769\\
[cd,1]=+1.09773\\
[ce,1]=-0.05852\\
[cf,1]=+0.26054\\
[dd,1]=+11.12064\\
[de,1]=-0.50528\\
[df,1]=-0.18790\\
[ee,1]=+2.26185\\
[ef,1]=-0.37202\\
[ff,1]=+5.61905\\[3pt]
\text{Und daraus ebenso}\\[3pt]
[nn,2]=+117763\\
[cn,2]=-115.81\\
[dn,2]=-153.95\\
[en,2]=+84.57\\
[fn,2]=-39.08\\
[cc,2]=+0.71612\\
[cd,2]=+1.11063\\
[ce,2]=-0.06392\\
[cf,2]=+0.25868\\
[dd,2]=+11.01463\\
[de,2]=-0.46088
\end{gathered}&
\begin{gathered}{}
[df,2]=-0.17265\\
[ee,2]=+2.24325\\
[ef,2]=-0.37841\\
[ff,2]=+5.61686\\[3pt]
\text{Daraus ferner}\\[3pt]
[nn,3]=+99034\\
[dn,3]=+25.66\\
[en,3]=+74.23\\
[fn,3]=+2.75\\
[dd,3]=+9.29213\\
[de,3]=-0.36175\\
[df,3]=-0.57384\\
[ee,3]=+2.23754\\
[ef,3]=-0.35532\\
[ff,3]=+5.52342\\[3pt]
\text{Daraus auf dieselbe Weise}\\[3pt]
[nn,4]=+98963\\
[en,4]=+75.23\\
[fn,4]=+4.33\\
[ee,4]=+2.22346\\
[ef,4]=-0.37766\\
[ff,4]=+5.48798\\[3pt]
\text{Daraus}\\[3pt]
[nn,5]=+96418\\
[fn,5]=+17.11\\
[ff,5]=+5.42383\\[3pt]
\text{Und daraus schließlich}\\[3pt]
[nn,6]=+96364
\end{gathered}
\end{array}
\]
}
Zur Bestimmung der Unbekannten haben wir also die folgenden sechs Gleichungen.
\[
\begin{aligned}
0={}&+17''11+5.42383\mathrm di\\
0={}&+75''23+2.22346\mathrm d\AscendingNode-0.37766\mathrm di\\
0={}&+25''66+9.29213\mathrm d\varphi-0.36175\mathrm d\AscendingNode-0.57384\mathrm di\\
0={}&-115''81+0.71612\mathrm d\Pi+1.11063\mathrm d\varphi-0.06392\mathrm d\AscendingNode+0.25868\mathrm di\\
0={}&-13854''+2458225\mathrm d\MeanMotionGlyph+62.13\mathrm d\Pi-510.58\mathrm d\varphi+213{,}84\mathrm d\AscendingNode\\
&+73.45\mathrm di\\[4pt]
0={}&-371''09+5.91569\mathrm dL+7203.91\mathrm d\MeanMotionGlyph-0.09344\mathrm d\Pi-2.28516\mathrm d\varphi\\
&-0.34664\mathrm d\AscendingNode-0.18194\mathrm di
\end{aligned}
\]
\SourceLeaf{29}
aus denen sich ergibt
\[
\begin{aligned}
\mathrm di&=-3''15\\
\mathrm d\AscendingNode&=-34''37\\
\mathrm d\varphi&=-4''29\\
\mathrm d\Pi&=+166''44\\
\mathrm d\MeanMotionGlyph&=+0''054335\\
\mathrm dL&=-3''06
\end{aligned}
\]

Die korrigierten elliptischen Elemente, die allen sechs Oppositionen möglichst genau genügen, sind somit die folgenden.
\begin{center}\fontsize{9}{11}\selectfont\setlength{\tabcolsep}{3pt}\begin{tabular}{lr}
Epoche der mittleren Länge 1803, am Göttinger Meridian&$221^0\,34'\,53''64$\\
Mittlere tropische tägliche Bewegung $\ .\ .\ .\ .\ .\ .\ .\ .\ .\ .\ .\ .$&$770''5010$\\
Länge des Perihels 1803 $\ .\ .\ .\ .\ .\ .\ .\ .\ .\ .\ .\ .\ .$&$121^0\,8'\,8''54$\\
Länge des aufsteigenden Knotens 1803 $\ .\ .\ .\ .\ .\ .\ .\ .\ .\ .$&$172\quad28\quad12.43$\\
Bahnneigung $\ .\ .\ .\ .\ .\ .\ .\ .\ .\ .\ .\ .\ .\ .\ .$&$34\quad37\quad28.35$\\
Exzentrizität $=\sin14^0\,9'\,59''79\ .\ .\ .\ .\ .\ .\ .\ .\ .\ .\ .$&$0.2447424$\\
Logarithmus der großen Halbachse $\ .\ .\ .\ .\ .\ .\ .\ .\ .\ .\ .\ .$&$0.4422071$
\end{tabular}\end{center}
\Article{15}
Setzen wir die eben gefundenen Werte der Korrekturen $\mathrm dL$, $\mathrm d\MeanMotionGlyph$ usw. in die zwölf Gleichungen aus Artikel 12 ein, so erhalten wir die folgenden Differenzen zwischen den beobachteten und berechneten Werten der heliozentrischen Längen und der geozentrischen Breiten.
\begin{center}\begin{tabular}{c|r|r}
Bei der Opposition&\multicolumn{2}{c}{Differenz}\\
des Jahres&\multicolumn{1}{c|}{der Länge}&\multicolumn{1}{c}{der Breite}\\\hline
1803&$-111''00$&$-8''31$\\
1804&$+59.18$&$-36.67$\\
1805&$+19.92$&$+0.07$\\
1807&$+85.77$&$+25.01$\\
1808&$+135.88$&$+28.72$\\
1809&$-216.54$&$+83.01$
\end{tabular}\end{center}
